% Chapter 2: Data Preprocessing
\chapter{Data Preprocessing}

\section{Overview}
This chapter describes the preprocessing pipeline applied to each dataset before integration. The preprocessing decisions were guided by data quality issues identified during EDA and requirements of the machine learning algorithms.

\section{Fire Dataset Preprocessing}

\subsection*{Data Cleaning}
\begin{itemize}
    \item \textbf{Missing coordinates:} Records excluded (spatial location essential)
    \item \textbf{Missing confidence:} Imputed with mode value
    \item \textbf{Missing FRP:} Imputed with median grouped by satellite sensor
    \item \textbf{Decision rationale:} Coordinate imputation would introduce spatial bias; other attributes can be reasonably estimated from similar records
\end{itemize}

\subsection*{Outlier Treatment}
\begin{itemize}
    \item FRP values capped at 99th percentile (extreme values from sensor errors)
    \item Points outside Algeria/Tunisia boundaries removed
\end{itemize}

\subsection*{Feature Extraction}
Temporal features extracted from acquisition date: Month, Season, Day/Night indicator. All timestamps standardized to UTC.

\section{Climate Dataset Preprocessing}

\subsection*{Handling Missing Data}
\begin{itemize}
    \item \textbf{Spatial gaps:} Bilinear interpolation from neighboring pixels
    \item \textbf{Temporal gaps:} Linear interpolation between adjacent months
    \item \textbf{Decision:} Climate variables are spatially and temporally autocorrelated, making interpolation valid
\end{itemize}

\subsection*{Normalization}
Temperature converted to Celsius; precipitation standardized to mm/month. Z-score normalization applied for modeling:
\begin{equation}
    z = \frac{x - \mu}{\sigma}
\end{equation}

\subsection*{Derived Features}
We created aggregated climate indices based on domain knowledge:
\begin{itemize}
    \item \textbf{Drought Index:} Combines precipitation deficit with temperature excess
    \item \textbf{Summer Heat Index:} Cumulative temperature above threshold during June-August
    \item \textbf{Decision:} These compound indices capture fire-relevant climate conditions better than raw variables
\end{itemize}

\section{Elevation Dataset Preprocessing}

\subsection*{Data Quality}
The GMTED2010 DEM required handling of NoData values at ocean and boundary areas through masking.

\subsection*{Derived Variables}
Topographic features calculated from elevation:
\begin{itemize}
    \item \textbf{Slope:} Using Horn's algorithm---steeper slopes affect fire spread
    \item \textbf{Terrain Complexity:} Local elevation standard deviation---rugged terrain affects fire behavior
    \item \textbf{Coastal Proximity:} Distance to coastline in km---coastal regions have distinct fire patterns
\end{itemize}

\section{Soil Dataset Preprocessing}

\subsection*{Missing Data}
Missing soil properties imputed using soil class averages from the HWSD database. Records with critical missing information excluded.

\subsection*{Normalization}
\begin{itemize}
    \item Organic carbon: Log-transformation followed by z-score normalization (skewed distribution)
    \item Texture fractions: Already percentages, no transformation needed
    \item pH: Direct z-score normalization (bounded scale)
\end{itemize}

\section{Data Integration}

\subsection*{Spatial Alignment}
All datasets aligned to WGS84 (EPSG:4326). Point-to-raster extraction used bilinear interpolation for continuous variables and nearest neighbor for categorical data.

\subsection*{Merging Strategy}
The integration followed a fire-centric approach:
\begin{enumerate}
    \item Fire point dataset as base
    \item Extract climate values at fire locations for corresponding months
    \item Extract elevation and topographic values at fire locations
    \item Spatial join with soil and land cover data
    \item Generate non-fire (negative) samples through stratified random sampling
\end{enumerate}

\begin{figure}[H]
    \centering
    \fbox{\parbox{0.85\textwidth}{\centering\vspace{3cm}\textbf{[PLACEHOLDER: Data Integration Flowchart]}\vspace{3cm}}}
    \caption{Data integration and merging workflow}
    \label{fig:merging_flowchart}
\end{figure}

\subsection*{Negative Sample Generation}
\textbf{Decision:} To create a classification dataset, we generated non-fire locations through stratified random sampling across the study area, matched temporally to fire observations to maintain climate consistency. This approach avoids spatial bias that could occur from simple random sampling.

\section{Final Dataset}

The merged dataset structure:

\begin{table}[H]
\centering
\caption{Final Merged Dataset Structure}
\label{tab:merged_structure}
\begin{tabular}{|l|l|c|}
\hline
\textbf{Category} & \textbf{Variables} & \textbf{Count} \\
\hline
Location & latitude, longitude & 2 \\
Climate & precipitation, temperature, drought\_index & 15 \\
Topography & elevation, slope, coastal\_proximity, terrain\_complexity & 6 \\
Soil & organic\_carbon, nitrogen, pH, texture & 8 \\
Target & fire (binary: 0/1) & 1 \\
\hline
\textbf{Total} & & \textbf{32+} \\
\hline
\end{tabular}
\end{table}

\subsection*{Train-Test Split}
Data split 80\%/20\% for training/testing with stratification on the target variable to maintain class proportions. Detailed preprocessing statistics are provided in Appendix~\ref{app:preprocessing}.
